\section{Cote 123: apposition, opposition, and rigid quadruples in type
A}\label{cote-123-apposition-opposition-and-rigid-quadruples-in-type-a}

\textbf{Scope and status.} These results answer the relevant original
questions for the complex simple adjoint groups
\(G=\operatorname{PGL}_n(\mathbf C)\), \(n\geq2\). All subgroup
automorphisms are complex algebraic automorphisms. No RH or
Dirichlet-GRH input is used. The arithmetic assertion at the end
concerns subgroup schemes individually defined over a specified subfield
\(k\subset\mathbf C\). Literature cutoff: 20 November 2026; source
inspection and proofs: 21 November 2026 UTC. No historical novelty is
asserted.

Grothendieck's letter of 22 October 1969, cote 123, physical PDF page
22, printed footer 21, asks for the classification of quadruples
\((T,B,T',B')\) with the maximal tori in Kostant apposition, whether the
two Borels are opposite for some or all choices, whether a distinguished
conjugacy class can be invariant under outer automorphisms, and what
automorphism groups these quadruples have. The original page was
visually inspected. The distinctions between these questions are
preserved below.

\begin{longtable}[]{@{}
  >{\raggedright\arraybackslash}p{(\columnwidth - 2\tabcolsep) * \real{0.5000}}
  >{\raggedright\arraybackslash}p{(\columnwidth - 2\tabcolsep) * \real{0.5000}}@{}}
\toprule\noalign{}
\begin{minipage}[b]{\linewidth}\raggedright
Result
\end{minipage} & \begin{minipage}[b]{\linewidth}\raggedright
Precise conclusion in type \(A_{n-1}\)
\end{minipage} \\
\midrule\noalign{}
\endhead
\bottomrule\noalign{}
\endlastfoot
CC-AP01 & Every ordered apposition pair is conjugate to the diagonal and
Fourier pair. \\
CC-AP02 & Opposition is exactly a specified list of nonvanishing Fourier
minors. \\
CC-AP03 & Every pair admits opposite Borel choices. All choices are
opposite if and only if \(n\) is prime. For composite \(n\), an explicit
common unipotent subgroup disproves opposition. \\
CC-AP04 & The inner-conjugacy classes are the explicit free quotient
\(K_n\backslash(S_n\times S_n)\), where
\(K_n=(\mathbf Z/n)^2\rtimes(\mathbf Z/n)^\times\). Their number is
\((n!)^2/(n^2\varphi(n))\). \\
CC-AP05 & For every odd \(n\geq3\), no such class is invariant under the
nontrivial outer automorphism; all quadruples then have trivial full
automorphism group. For even \(n\geq4\), outer-invariant examples exist,
with automorphism group \(C_2\). An exact finite criterion covers every
class. \\
CC-AP06 & A quadruple whose four subgroup schemes are defined over
\(k\subset\mathbf C\) exists in split \(\operatorname{PGL}_{n,k}\) if
and only if \(k\) contains a primitive \(n\)-th root of unity. In
particular there is no such quadruple over \(\mathbf Q\) for
\(n\geq3\). \\
CC-AP07 & Adding a generating principal element of \(T\cap N_G(T')\)
gives the free quotient \((\mathbf Z/n)^2\backslash(S_n\times S_n)\),
with \((n!)^2/n^2\) classes. \\
\end{longtable}

\subsection{Inspected primary
convention}\label{inspected-primary-convention}

Use Kostant's convention of a Cartan \(\mathfrak t'\) in apposition to
\(\mathfrak t\) with respect to a principal element \(P\). The intrinsic
characterization needed here is

\[
 \mathfrak t=\mathfrak g^{\operatorname{Ad}P},
 \qquad
 \operatorname{Ad}P(\mathfrak t')=\mathfrak t',
 \qquad
 \operatorname{Ad}P|_{\mathfrak t'}\text{ is a Coxeter element}.
 \tag{CC-AP.1}
\]

A lift of a Weyl-group Coxeter element produces this configuration, and
apposition has this Coxeter action. These are Kostant (1959),
\textbf{Theorem 8.6 and the immediately following Corollary 8.6,
p.~1026}. Their statements were read in the OCR of the author's
collected-paper reprint. Kostant's own later account provides an
independently accessible primary check: \textbf{Theorem 0.1 and the
following apposition sentence, arXiv:1003.0046v1 p.~2; Theorems 1.2--1.3
p.~6; Proposition 1.8 and Theorem 1.9 p.~8}. His principal element is
\(\exp(2\pi i\rho^\vee/h)\); regularity alone is not its definition.

The original definition in 1959, §7.3, p.~1018, was located through
later references but was \textbf{not directly inspected} in the original
scan. No claim of such inspection is made. The proof below uses the
inspected equivalent characterization (CC-AP.1), together with the
cyclic-element construction explicitly identified as apposition by
Kostant himself in 2010. The short proof of his Theorem 1.9 was read;
the full 1959 structure theorem is a cited input.

When using (CC-AP.1) at group level, take the connected centralizer
determined by its Lie algebra. In an adjoint group the full centralizer
of a regular semisimple group element can be disconnected.

\subsection{CC-AP01 --- The diagonal and Fourier normal
form}\label{cc-ap01-the-diagonal-and-fourier-normal-form}

Put \(I=\mathbf Z/n\mathbf Z\), fix \(\zeta=e^{2\pi i/n}\), and use the
basis \(e_0,\ldots,e_{n-1}\) of \(V=\mathbf C^n\). Define

\[
 C e_j=e_{j+1},\qquad
 Qe_j=\zeta^j e_j,\qquad
 v_a=\sum_{j\in I}\zeta^{ja}e_j,
 \qquad F=(\zeta^{ja})_{j,a\in I}.
\]

Let \(T_0\) be the diagonal maximal torus of \(\operatorname{PGL}(V)\),
and set \(T_1=FT_0F^{-1}\).

\textbf{Proposition.} The ordered pair \((T_0,T_1)\) is in Kostant
apposition, and every ordered apposition pair in
\(\operatorname{PGL}_n(\mathbf C)\) is conjugate to it.

\textbf{Proof.} The vectors \(v_a\) form a basis because the usual
Vandermonde determinant of the distinct numbers \(\zeta^a\) is nonzero.
Direct calculation gives

\[
 Cv_a=\zeta^{-a}v_a,\qquad Qv_a=v_{a+1},
 \qquad QCQ^{-1}=\zeta C.
\]

Thus \(T_1\) is the connected centralizer of \(C\). The adjoint action
of \(Q\) fixes exactly the diagonal Cartan and acts on the \(n\)
eigenlines of \(T_1\) by an \(n\)-cycle. Such a cycle is a Coxeter
element of \(S_n\). This verifies apposition by (CC-AP.1). Equivalently,
\(C=\sum_{j=0}^{n-2}E_{j+1,j}+E_{0,n-1}\) is precisely the sum of the
simple-root vectors and the negative-highest-root vector for the reverse
standard positive system. All coefficients are nonzero, so Kostant's
cyclic-element construction applies. Projectively
\(Q=\exp(2\pi i\rho^\vee/n)\) for this positive system.

Conversely, take a pair \((T,T')\) and a principal \(P\) as in
(CC-AP.1). Conjugate \(T'\) to \(T_0\). The action of \(P\) on its
eigenlines is a Coxeter element, hence an \(n\)-cycle. After a
permutation of the eigenlines, a matrix representative of \(P\) has the
form

\[
 A e_j=c_j e_{j+1},\qquad c_j\ne0.
\]

Choose \(c\) with \(c^n=\prod_jc_j\) and choose nonzero diagonal entries
\(d_j\) recursively by \(d_{j+1}=c\,d_j/c_j\). The last equation agrees
with the first by the displayed product relation. Conjugation by
\(\operatorname{diag}(d_j)\) changes \(A\) to \(cC\), which represents
\(C\) in \(\operatorname{PGL}(V)\). Therefore the ordered pair becomes
\((T_1,T_0)\), since the fixed Lie algebra of \(\operatorname{Ad}C\) is
the Cartan centralizing \(C\).

Finally \(F^2=nR\), where \(Re_j=e_{-j}\). The permutation \(R\)
normalizes \(T_0\). It follows that \(F^{-1}T_0F=T_1\), while
\(F^{-1}T_1F=T_0\). Conjugating by \(F^{-1}\) therefore exchanges the
two tori and gives the asserted ordered normal form. \(\square\)

\subsection{CC-AP02 --- The exact opposition
test}\label{cc-ap02-the-exact-opposition-test}

Write permutations \(\pi,\sigma\in S_I\) as ordered lists. Every Borel
containing \(T_0\) stabilizes exactly one of the flags

\[
 {\cal E}^{\pi}_k=\langle e_{\pi(0)},\ldots,e_{\pi(k-1)}\rangle ,
\]

and every Borel containing \(T_1\) is obtained from

\[
 {\cal V}^{\sigma}_r=\langle v_{\sigma(0)},\ldots,v_{\sigma(r-1)}\rangle .
\]

Here \(0\leq k,r\leq n\). This is the full-flag description of Borels in
\(\operatorname{PGL}(V)\): the invariant subspaces for a maximal
diagonal torus are sums of its distinct one-dimensional weight spaces.

\textbf{Proposition.} The Borels
\(B_\pi=\operatorname{Stab}({\cal E}^\pi)\) and
\(B'_\sigma=\operatorname{Stab}({\cal V}^\sigma)\) are opposite if and
only if, for every \(1\leq k\leq n-1\),

\[
 \det\bigl(\zeta^{ja}\bigr)_{
 j\in\{\pi(k),\ldots,\pi(n-1)\},
 \ a\in\{\sigma(0),\ldots,\sigma(n-k-1)\}}\ne0.
 \tag{CC-AP.2}
\]

The order chosen on either displayed set affects only the sign.

\textbf{Proof.} First recall directly why opposite flags are
characterized by

\[
 {\cal E}_k\cap{\cal V}_{n-k}=0
 \quad(1\leq k\leq n-1).
 \tag{CC-AP.3}
\]

Two full flags admit a simultaneous adapted basis. One construction uses
the dimensions \(r_{ij}=\dim({\cal E}_i\cap{\cal V}_j)\): the
nonnegative integers

\[
 r_{ij}-r_{i-1,j}-r_{i,j-1}+r_{i-1,j-1}
\]

are the dimensions of the corresponding successive quotients of the
intersections. Their row and column sums are one, so the nonzero
positions form a permutation matrix. Choosing a vector at each of these
positions gives the simultaneous basis, by induction on the
intersections.

Order that basis according to the first flag, and let \(\eta(i)\) be the
position of its \(i\)-th vector in the second flag. If an increasing
pair \(i<j\), \(\eta(i)<\eta(j)\), occurs, the nontrivial root subgroup
sending the \(j\)-th basis vector to itself plus a multiple of the
\(i\)-th preserves both flags. Hence the common stabilizer contains a
nontrivial unipotent subgroup and is not a torus. If there is no
increasing pair, \(\eta\) is the reverse ordering. The common stabilizer
is then exactly the diagonal torus in this basis. The reverse ordering
is also equivalent to (CC-AP.3). This proves the equivalence with
opposition in the sense used in the letter.

For the flags in the statement, the quotient map
\(V\to V/{\cal E}^{\pi}_k\), restricted to \({\cal V}^{\sigma}_{n-k}\),
is represented by the square submatrix in (CC-AP.2). Its domain and
codomain both have dimension \(n-k\). Its kernel is the intersection in
(CC-AP.3). Thus (CC-AP.2) and (CC-AP.3) are equivalent. \(\square\)

\subsection{CC-AP03 --- ``Some'' always; ``all'' exactly in prime
dimension}\label{cc-ap03-some-always-all-exactly-in-prime-dimension}

\textbf{Theorem.} For \(G=\operatorname{PGL}_n(\mathbf C)\):

\begin{enumerate}
\def\labelenumi{\arabic{enumi}.}
\tightlist
\item
  Every apposition pair admits opposite containing Borels. More
  strongly, after the normal form of CC-AP01, the Fourier flag in the
  ordering \(0,1,\ldots,n-1\) is opposite to every coordinate flag, and
  the coordinate flag in that ordering is opposite to every Fourier
  flag.
\item
  Every quadruple with tori in apposition has opposite Borels if and
  only if \(n\) is prime.
\item
  If \(n\) is composite, the failure can be witnessed by a rank-one
  nilpotent matrix whose associated one-parameter unipotent group lies
  in both Borels.
\end{enumerate}

\textbf{Proof.} For the first assertion, fix the Fourier ordering
\(\sigma(a)=a\). For any coordinate ordering and any \(k\), the minor in
(CC-AP.2) has columns with exponents \(0,\ldots,n-k-1\) and distinct row
parameters \(\zeta^j\). It is a Vandermonde determinant and is nonzero.
Conversely, for the standard coordinate ordering its rows have
consecutive exponents \(k,\ldots,n-1\). Factoring \(\zeta^{ka}\) from
the column indexed by \(a\) again leaves a Vandermonde determinant in
the distinct column parameters \(\zeta^a\). This proves both
strengthened assertions.

If \(n=p\) is prime, every square Fourier minor is nonzero by the
Chebotarev minor theorem, in the exact primary formulation \textbf{Tao
(2005), Lemma 1.3, p.~123, proof pp.~124--125}; equivalently
arXiv:math/0308286v6, Lemma 1.3, pp.~3--4. Applying that theorem to
(CC-AP.2) proves opposition for every pair of containing Borels. CC-AP01
then covers every apposition pair. The cited minor theorem is a known
input; this application does not claim a new proof of it.

For the converse write \(n=ab\), where \(a,b>1\). The submatrix with
rows \(\{0,a\}\) and columns \(\{0,b\}\) is

\[
 \begin{pmatrix}1&1\\1&\zeta^{ab}\end{pmatrix}
 =\begin{pmatrix}1&1\\1&1\end{pmatrix}.
\]

Choose the coordinate ordering so that its last two indices are \(0,a\),
and choose the Fourier ordering so that its first two indices are
\(0,b\). Equation (CC-AP.2) fails at \(k=n-2\).

For a direct group-theoretic witness put

\[
 x=v_0-v_b,\qquad
 \lambda=e_0^*-e_a^*,\qquad
 N=x\otimes\lambda.
 \tag{CC-AP.4}
\]

The vector \(x\ne0\) belongs both to \({\cal E}_{n-2}^{\pi}\) and to
\({\cal V}_2^\sigma\), because its coordinates at \(0,a\) vanish. The
functional \(\lambda\) vanishes on both these subspaces: on the
coordinate subspace this is immediate, and on \(v_0,v_b\) it follows
from \(\zeta^{ab}=1\). Hence \(N\ne0\), \(N^2=0\), and \(I+tN\) is a
nontrivial one-parameter unipotent subgroup.

Every earlier step of the coordinate flag is contained in
\(\ker\lambda\), and every later step contains \(x\). The same holds for
the Fourier flag, with its dividing step at dimension two. Therefore
\(I+tN\) preserves both complete flags. Their common stabilizer contains
this nontrivial unipotent subgroup and cannot be a torus.

In the smallest example \(n=4\), take \(\zeta=i\), coordinate ordering
\((1,3,0,2)\), and Fourier ordering \((0,2,1,3)\). Rescaling (CC-AP.4)
gives the particularly simple witness

\[
 N=(e_1+e_3)(e_0^*-e_2^*).
\]

This is a counterexample inside the original complex simple adjoint
group of type \(A_3\), with apposition established before the flag
calculation. \(\square\)

\subsection{CC-AP04 --- Explicit classification and inner
rigidity}\label{cc-ap04-explicit-classification-and-inner-rigidity}

Let

\[
 K_n=N_G(T_0)\cap N_G(T_1),\qquad
 \Omega_n=S_I\times S_I.
\]

For \(a\in I^\times\) and \(b,d\in I\), define the projective class of
the matrix

\[
 g_{a,b,d}e_j=\zeta^{ad j}e_{aj+b}.
 \tag{CC-AP.5}
\]

\textbf{Theorem.} These are exactly the elements of \(K_n\). On the
coordinate and Fourier eigenline labels their action is

\[
 j\longmapsto aj+b,\qquad
 c\longmapsto a^{-1}c+d.
 \tag{CC-AP.6}
\]

Consequently

\[
 K_n\simeq(I\times I)\rtimes I^\times,\qquad
 a\cdot(b,d)=(ab,a^{-1}d).
\]

The inner-conjugacy classes of ordered quadruples are exactly

\[
 K_n\backslash\Omega_n,\qquad
 (a,b,d)\cdot(\pi,\sigma)
 =(a\pi+b,\ a^{-1}\sigma+d).
 \tag{CC-AP.7}
\]

This action is free. Thus every quadruple has trivial inner stabilizer,
and the number of classes is

\[
 \frac{(n!)^2}{n^2\varphi(n)}.
 \tag{CC-AP.8}
\]

\textbf{Proof.} An element normalizing \(T_0\) is monomial in the
coordinate basis, and an element normalizing \(T_1\) is monomial in the
Fourier basis. Thus a representative \(g\) of an element of \(K_n\)
satisfies, for permutations \(\pi_0,\sigma_0\) of \(I\) and nonzero
scalars \(u_j,w_c\),

\[
 ge_j=u_j e_{\pi_0(j)},\qquad
 gv_c=w_c v_{\sigma_0(c)}.
\]

Comparison of the coordinate indexed by \(\pi_0(j)\) yields

\[
 u_j\zeta^{jc}=w_c\zeta^{\pi_0(j)\sigma_0(c)}.
\]

Divide this equation by its instance at \(c=0\), and then use its
instance at \(j=0\) to remove the remaining scalar. The result is the
congruence

\[
 jc=(\pi_0(j)-\pi_0(0))(\sigma_0(c)-\sigma_0(0))
 \pmod n.
 \tag{CC-AP.9}
\]

At \(j=c=1\), the two factors on the right are units and are inverses.
Put \(a=\pi_0(1)-\pi_0(0)\), \(b=\pi_0(0)\), and \(d=\sigma_0(0)\).
Taking \(c=1\) and then \(j=1\) in (CC-AP.9) gives

\[
 \pi_0(j)=aj+b,\qquad \sigma_0(c)=a^{-1}c+d.
\]

The equation at \(c=0\) then gives \(u_j=u\,\zeta^{ad j}\), with a
scalar \(u\) independent of \(j\). This proves that the projective
matrix is (CC-AP.5).

Conversely (CC-AP.5) gives

\[
 g_{a,b,d}v_c
 =\zeta^{-b(a^{-1}c+d)}v_{a^{-1}c+d},
\]

so all the displayed matrices normalize both tori. Their action on the
two sets of eigenlines is faithful. Composition of the affine
permutations gives

\[
 (a,b,d)(a',b',d')
 =(aa',\,ab'+b,\,a^{-1}d'+d),
\]

which proves the claimed semidirect-product structure and its order
\(n^2\varphi(n)\).

By CC-AP01 every quadruple has a representative with torus pair
\((T_0,T_1)\), and precisely \(K_n\) remains available for conjugating
such representatives. This proves the orbit classification (CC-AP.7). If
a \(K_n\)-element preserves a full flag made from eigenlines, its
permutation of those eigenlines preserves every initial segment and is
therefore the identity. Preserving both flags makes both affine
permutations in (CC-AP.6) the identity. It follows that \(a=1,b=d=0\),
and the projective transformation is the identity. The action is free,
so counting proves (CC-AP.8). This also proves inner rigidity for every
original quadruple, not merely for an opposite one. \(\square\)

For example, there is one class in type \(A_1\), two in type \(A_2\),
and eighteen in type \(A_3\). Formula (CC-AP.7) is an explicit finite
classification in type A; it is not a classification in the other simple
types.

\subsection{CC-AP05 --- Outer action, distinguished classes, and full
automorphisms}\label{cc-ap05-outer-action-distinguished-classes-and-full-automorphisms}

Let \(\theta(g)=(g^{-1})^{\mathsf T}\). For \(n\geq3\) it represents the
nontrivial element of the outer automorphism group of \(G\).

\textbf{Theorem.} On the quotient (CC-AP.7), \(\theta\) acts by

\[
 (\pi,\sigma)\longmapsto
 \bigl(\pi^{\mathrm{rev}},-\sigma^{\mathrm{rev}}\bigr),
 \quad
 \pi^{\mathrm{rev}}(i)=\pi(n-1-i).
 \tag{CC-AP.10}
\]

A class is fixed if and only if some \(a\in I^\times\), \(b,d\in I\)
satisfy, for every \(0\leq i<n\),

\[
 a\pi(i)+b=\pi(n-1-i),\qquad
 a^{-1}\sigma(i)+d=-\sigma(n-1-i).
 \tag{CC-AP.11}
\]

For \(n\geq3\), the full automorphism group of an individual quadruple
is trivial if this condition fails and is \(C_2\) if it holds. If \(n\)
is odd, the condition always fails. If \(n\geq4\) is even, it holds for
explicit examples.

\textbf{Proof.} We first justify the outer-automorphism statement in the
precise algebraic category being used. Any algebraic automorphism of
\(G\), after an inner adjustment, preserves a chosen Borel and its
maximal torus. Its action on the roots then preserves the positive
system and induces an automorphism of the chain \(A_{n-1}\). There are
only identity and reversal, with only identity when \(n=2\). The
composition of \(\theta\) with conjugation by the reverse permutation
matrix preserves this Borel and torus and induces reversal.

If the induced root action is the identity, the automorphism is the
identity on the torus, since its character lattice for the adjoint group
is generated by the roots. On the Lie algebra root spaces it has the
form \(E_{ij}\mapsto c_{ij}E_{ij}\). Preservation of brackets gives
\(c_{ij}c_{ji}=1\), and, for distinct \(i,j,k\),
\(c_{ik}=c_{ij}c_{jk}\). Fixing one index shows \(c_{ij}=s_i/s_j\) for
suitable \(s_i\ne0\); the two-index case follows directly from
\(c_{ij}c_{ji}=1\). Thus conjugation by \(\operatorname{diag}(s_i)\)
gives the same differential. Automorphisms of the connected complex
group are determined by their differential, so the automorphism is
inner. Finally, for \(n\geq3\), reversal of the simple roots cannot
arise from an inner automorphism preserving this Borel and torus: its
conjugating element would belong to \(B\cap N_G(T)=T\), whose action on
the roots is the identity. This proves \(\operatorname{Out}(G)=C_2\)
with representative \(\theta\) when \(n\geq3\), and
\(\operatorname{Out}(G)=1\) when \(n=2\).

Both \(T_0\) and \(T_1\) are preserved by \(\theta\): inverse transpose
preserves diagonal matrices and sends the algebra of circulant matrices
to itself. On a full flag \(\theta\) sends the stabilizer to the
stabilizer of the flag of annihilators, in reverse order, using the
standard bilinear form to identify \(V^*\) with \(V\). The dual
coordinate basis is the same coordinate basis. For the Fourier basis one
has

\[
 v_a^{\mathsf T}v_b=n\,\mathbf 1_{a+b=0},
\]

so its dual basis has labels \(-a\). This proves (CC-AP.10). Combining
it with (CC-AP.7) proves the equivalence with (CC-AP.11).

By CC-AP04 the inner stabilizer is trivial. Consequently the
automorphism group of a quadruple injects into
\(\operatorname{Out}(G)\), and its image is exactly the stabilizer of
its inner-conjugacy class. In particular, if (CC-AP.11) holds, an inner
adjustment of \(\theta\) fixes the quadruple. Its square is inner and
fixes the quadruple, hence is the identity. This proves the asserted
automorphism-group classification.

Now suppose \(n\geq3\) is odd and (CC-AP.11) holds. The affine
permutation

\[
 R(j)=aj+b
\]

reverses an ordering of \(n\) elements. It is an involution with exactly
one fixed point. Therefore \(a^2=1\pmod n\), and the equation
\((a-1)j=-b\) has exactly one solution. The latter fact implies
\(\gcd(a-1,n)=1\). Multiplying the identity \((a-1)(a+1)=0\) by the
inverse of \(a-1\) gives \(a=-1\pmod n\).

The second equation of (CC-AP.11) says that

\[
 S(c)=-a^{-1}c-d
\]

also reverses an ordering of \(n\) elements. With \(a=-1\), this is the
translation \(S(c)=c-d\). Since it is an involution and \(n\) is odd,
\(2d=0\) implies \(d=0\). Then \(S\) is the identity, which cannot
reverse an ordering of at least three distinct elements. This
contradiction proves the assertion for every odd \(n\geq3\), including
composite \(n\).

For an even example write \(n=2m\), with \(m\geq2\), and take

\[
 \pi=(0,1,\ldots,m-1,\ 2m-1,2m-2,\ldots,m),\qquad
 \sigma=(0,1,\ldots,2m-1).
\]

Equation (CC-AP.11) holds with \(a=1,b=m,d=1\). The corresponding class
is outer-invariant, and its automorphism group is \(C_2\), as already
proved. \(\square\)

\textbf{Direct answer to the original distinguished-class question.} In
every type \(A_{n-1}\) with odd \(n\geq3\), there is no inner-conjugacy
class invariant under all outer automorphisms. Thus a distinguished
class satisfying Grothendieck's stated necessary invariance condition
cannot exist there. In type \(A_2\), the two classes of CC-AP04 are
exchanged. Every quadruple in these odd types has trivial full algebraic
automorphism group. In type \(A_1\) the sole class is invariant and its
automorphism group is trivial. For even \(n\geq4\) the theorem supplies
invariant classes and the exact automorphism test, but does not assert a
unique canonical choice among them.

For a completely explicit description in type \(A_2\), the two classes
are distinguished by
\(\operatorname{sgn}(\pi)\operatorname{sgn}(\sigma)\). Translations of
\(\mathbf Z/3\) are even permutations, while multiplication by \(-1\) is
odd on both sets of labels, so this product is \(K_3\)-invariant. The
count of two classes proves it is complete. Representatives are
\((\pi,\sigma)=((0,1,2),(0,1,2))\) and \(((0,1,2),(0,2,1))\). Reversal
changes the row sign, whereas negation followed by reversal leaves the
Fourier sign unchanged; thus \(\theta\) exchanges the two classes
directly. Both representatives have opposite Borels by CC-AP03.

\subsection{CC-AP06 --- The exact split-field
obstruction}\label{cc-ap06-the-exact-split-field-obstruction}

Let \(k\subset\mathbf C\) be a field, and retain \(n\geq2\). Here a
\(k\)-defined quadruple means that each of \(T,B,T',B'\) is a subgroup
scheme of the split group \(\operatorname{PGL}_{n,k}\) defined over
\(k\), and that after extending to \(\mathbf C\) it has the apposition
property above.

\textbf{Theorem.} Such a quadruple exists if and only if
\(\zeta_n\in k\) for a primitive \(n\)-th root of unity \(\zeta_n\).

\textbf{Proof.} A \(k\)-defined Borel in split \(\operatorname{PGL}_n\)
stabilizes a full flag over \(k\). A \(k\)-defined maximal torus inside
it is split. To see the latter assertion concretely, lift the torus and
Borel to \(\operatorname{GL}_n\) and choose a \(k\)-basis adapted to the
flag. The lifted torus has dimension \(n\) and is a torus after base
change. Projection of the upper triangular group onto its diagonal
quotient restricts to an injective map on this torus: its kernel would
be both a torus subgroup and a unipotent subgroup, and hence trivial in
characteristic zero. By equal dimension this map is an isomorphism onto
\((\mathbf G_m)^n\). Thus the torus is split, and its one-dimensional
weight spaces admit a basis over \(k\).

Apply this argument to both \(T\subset B\) and \(T'\subset B'\). Choose
a \(k\)-basis diagonalizing the first torus. There is an
\(A\in\operatorname{GL}_n(k)\) such that the second torus is
\(AT_0A^{-1}\). CC-AP01 gives a complex transformation preserving the
first torus and carrying the second to \(FT_0F^{-1}\). Transformations
preserving \(T_0\) are monomial, so this says exactly that \(A\) is
obtained from \(F\) by row and column permutations and nonzero row and
column scalings.

All entries of \(A\) are therefore nonzero. Its cross-ratios

\[
 \frac{A_{ic}A_{jd}}{A_{id}A_{jc}}\qquad(i\ne j,\ c\ne d)
 \tag{CC-AP.12}
\]

belong to \(k\). Row and column scalings leave these ratios unchanged,
and row and column permutations only relabel them. For the Fourier
matrix, taking rows \(0,1\) and columns \(0,1\) gives the ratio
\(\zeta\). Hence \(\zeta\in k\).

Conversely, if a primitive \(\zeta\) belongs to \(k\), the Fourier
matrix, both tori, and the coordinate and Fourier flags are all defined
over \(k\). They give the required quadruple by CC-AP01. \(\square\)

For \(n=3\) and \(k=\mathbf Q\), this is an independently proved
arithmetic nonsplitness obstruction: a full rational apposition
quadruple does not exist in split \(\operatorname{PGL}_3\), even though
the diagonal torus and the connected centralizer of the rational
cyclic-shift matrix give a rationally defined pair of tori. The second
torus in that pair is not split over \(\mathbf Q\).

The supplied SA20 intake states an arithmetic nonsplitness result
labelled D.016 for ``RigKos in \(A_2/\mathbf Q\)'' but explicitly says
that Part D was read only at statement level. The full definition and
proof of that object were not supplied here. CC-AP06 therefore
establishes the concrete subgroup-scheme assertion above and does not
identify it with an unseen moduli functor, a coarse moduli point, or
every possible twisted form.

\subsection{CC-AP07 --- Classification with a specified generating
principal
element}\label{cc-ap07-classification-with-a-specified-generating-principal-element}

The original letter also asks for the classification when a generating
principal element of the finite group \(T\cap N_G(T')\) is included in
the data.

\textbf{Corollary.} In type \(A_{n-1}\), this group is cyclic of order
\(n\). The inner-conjugacy classes of the ordered data
\((T,B,T',B',P)\), where \(P\) generates that group, are

\[
 K_n\backslash(\Omega_n\times I^\times),
 \quad
 (a,b,d)\cdot(\pi,\sigma,u)
 =(a\pi+b,\ a^{-1}\sigma+d,\ a^{-1}u).
 \tag{CC-AP.13}
\]

Equivalently, after normalizing \(P=Q\), they are

\[
 (I\times I)\backslash\Omega_n,\qquad
 (b,d)\cdot(\pi,\sigma)=(\pi+b,\sigma+d).
 \tag{CC-AP.14}
\]

Both actions are free and their common number of orbits is
\((n!)^2/n^2\).

\textbf{Proof.} In the normal form, an element of \(T_0\cap N_G(T_1)\)
belongs to \(K_n\) and has trivial permutation of the coordinate
eigenlines. By (CC-AP.6), this forces \(a=1,b=0\), so the element is
\(g_{1,0,d}=Q^d\). Conversely every such power normalizes \(T_1\). Hence

\[
 T_0\cap N_G(T_1)=\langle Q\rangle\simeq C_n.
\]

Its generators are precisely \(Q^u\), \(u\in I^\times\). They are
principal: each is conjugate to \(Q\) by permuting its \(n\) distinct
eigenvalues. Direct conjugation, or (CC-AP.5), gives in the projective
group

\[
 g_{a,b,d}Q^u g_{a,b,d}^{-1}=Q^{a^{-1}u}.
\]

This proves (CC-AP.13). Its action is free because its action on the
first two coordinates is already free by CC-AP04. Therefore its orbit
count is

\[
 \frac{(n!)^2\varphi(n)}{n^2\varphi(n)}
 =\frac{(n!)^2}{n^2}.
\]

Given any generator exponent \(u\), choose \(a=u\) to change it to
\(1\). The remaining transformations fixing \(u=1\) have \(a=1\) and are
exactly the two independent translations in (CC-AP.14). This proves the
alternative classification, its freeness, and its count. \(\square\)

These two counts agree with the letter's \(|W|^2/(hz\varphi(h))\) and
\(|W|^2/(hz)\): in type \(A_{n-1}\), \(W=S_n\), \(h=n\), and \(z=n\) is
the order of the center of the simply connected cover
\(\operatorname{SL}_n\).

\subsection{Primary bibliography and inspection
record}\label{primary-bibliography-and-inspection-record}

Grothendieck, A. (1969, October 22). \emph{Letter to B. Kostant}.
Montpellier archives, cote 123, physical PDF page 22, footer 21.
https://grothendieck.umontpellier.fr/123.pdf. The original page was
visually read, including the distinct classification, opposition,
distinguished-class, and automorphism questions.

Kostant, B. (1959). The principal three-dimensional subgroup and the
Betti numbers of a complex simple Lie group. \emph{American Journal of
Mathematics, 81}(4), 973--1032. https://doi.org/10.2307/2372999. Theorem
8.6 and the following Corollary 8.6, p.~1026, were read in the OCR of
the author's reprint in \emph{Collected Papers}, Vol. I, pp.~130--189
(2009), https://doi.org/10.1007/b94535\_11. The source scan itself was
not obtained; original §7.3, p.~1018, is a located but uninspected
reference. OCR inspected at
https://dokumen.pub/collected-papers-volume-i-1955-1966-2009nbsped-9780387095820-0387095829.html.

Kostant, B. (2010). Experimental evidence for the occurrence of \(E_8\)
in nature and the radii of the Gosset circles. \emph{Selecta
Mathematica, 16}, 419--438. https://doi.org/10.1007/s00029-010-0023-1.
Inspected author preprint: https://arxiv.org/pdf/1003.0046v1 (1 March
2010). Exact locators in this version: Theorem 0.1 and apposition
sentence p.~2; Proposition 1.1 p.~5; Theorems 1.2--1.3 p.~6; Theorem 1.7
and principal-element discussion p.~7; Proposition 1.8 and Theorem 1.9,
including its complete short proof, p.~8.

Tao, T. (2005). An uncertainty principle for cyclic groups of prime
order. \emph{Mathematical Research Letters, 12}(1), 121--127.
https://doi.org/10.4310/MRL.2005.v12.n1.a11. Lemma 1.3 p.~123; proof
pp.~124--125; preliminary Lemma 1.2 p.~123. The journal PDF was directly
downloaded and the statements and complete indicated proofs read:
https://www.intlpress.com/site/pub/files/\_fulltext/journals/mrl/2005/0012/0001/MRL-2005-0012-0001-a011.pdf.
Also inspected: https://arxiv.org/pdf/math/0308286v6 (22 July 2004),
Lemmas 1.2--1.3 pp.~3--4. The prime Fourier-minor statement is credited
there to Chebotarev; no original Chebotarev theorem number is invented.

\textbf{Proof provenance.} CC-AP01--07 give complete proofs of the
stated type-A deductions, with Kostant's apposition/Coxeter
characterization and the prime Fourier-minor theorem explicitly retained
as known inputs. The finite normalizer and outer-action calculations,
the composite unipotent witness, the field-of-definition obstruction,
and the classification with a specified generator are proved here. This
is neither a proof of RH nor a classification of apposition quadruples
in every simple type.
